% from ega4-11-fr.tex
\oldpage[IV-3]{172}
\begin{proposition}[11.10.6]
\label{IV.11.10.6-fr}
Les notations étant celles de (11.10.5), supposons que $X$ soit un préschéma
sur un corps $k$, et qu'en posant $S=\operatorname{Spec}(k)$, on ait
$X'=X\times_SS'$, où $S'$ est un $k$-préschéma quelconque. Alors, si la famille
$(f_\lambda)$ est schématiquement dominante, il en est de même de
$(f'_\lambda)$.
\end{proposition}

\begin{corollary}[11.10.7]
\label{IV.11.10.7-fr}
Soient $X$ un préschéma sur un corps $k$, $(Z_\lambda)_{\lambda\in L}$ une
famille de $k$-préschémas géométriquement réduits sur $k$, et pour chaque
$\lambda\in L$, soit $f_\lambda:Z_\lambda\to X$ un $k$-morphisme. Désignons
par $Y$ le sous-préschéma réduit de $X$ ayant pour espace sous-jacent
l'adhérence de $\bigcup_{\lambda\in L}f_\lambda(Z_\lambda)$. Soient $k'$ une
extension de $k$, $X'$, $Z'_\lambda$,
$f'_\lambda:Z'_\lambda\to X'$ les préschémas et morphismes déduits de $X$,
$Z_\lambda$ et $f_\lambda$ par extension de la base à $k'$. Alors, si $Y'$ est
le sous-préschéma réduit de $X'$ ayant pour espace sous-jacent l'adhérence de
$\bigcup_{\lambda\in L}f'_\lambda(Z'_\lambda)$, on a
$Y'=Y\otimes_k k'$. En particulier, $Y$ est géométriquement réduit sur $k$.
\end{corollary}

\begin{proof}
Comme les $Z_\lambda$ sont réduits, les morphismes $f_\lambda$ se factorisent
en
\[
Z_\lambda\xrightarrow{g_\lambda}Y\xrightarrow{j}X,
\]
où $j$ est l'injection canonique ($\mathbf I$, 5.2.2). Il résulte alors de
(11.10.4) que $(g_\lambda)$ est une famille schématiquement dominante. Posons
$Y'_1=Y\otimes_k k'$, sous-préschéma fermé de $X'$, et soit $g'_\lambda$ le
morphisme déduit de $g_\lambda$ par extension de la base à $k'$, de sorte que
$f'_\lambda$ se factorise en
\[
Z'_\lambda\xrightarrow{g'_\lambda}Y'_1\xrightarrow{j'}X',
\]
où $j'$ est l'injection canonique. Il résulte de (11.10.6) que la famille
$(g'_\lambda)$ est schématiquement dominante. Mais par hypothèse les
$Z'_\lambda$ sont \emph{réduits}, donc ($\mathbf I$, 5.2.2) les
$g'_\lambda$ se factorisent en $Z'_\lambda\xrightarrow{h'_\lambda}Y'\to Y'_1$~; on conclut donc
de (11.10.2) que $Y'=Y'_1$ et $h'_\lambda=g'_\lambda$, ce qui établit le
corollaire.
\end{proof}

\begin{definition}[11.10.8]
\label{IV.11.10.8-fr}
Les notations étant celles de (11.10.5), supposons que $X$ soit un
$S$-préschéma. On dit que $(f_\lambda)$ est \emph{universellement
schématiquement dominante} (relativement à $S$) si, pour tout changement de
base $S'\to S$, la famille $(f'_\lambda)$ correspondant au changement de base
$X'=X\times_SS'\to X$, est schématiquement dominante.
\end{definition}

Lorsque $S$ est le spectre d'un corps, la prop. (11.10.6) signifie donc qu'une
famille schématiquement dominante est universellement schématiquement
dominante (relativement à $S$).

Lorsque les $f_\lambda$ sont des immersions canoniques $Z_\lambda\to X$, on
dira aussi que la famille des sous-préschémas $Z_\lambda$ est
\emph{universellement schématiquement dense dans $X$} (relativement à $S$) au
lieu de dire que la famille $(f_\lambda)$ est universellement schématiquement
dominante (relativement à $S$).

\begin{theorem}[11.10.9]
\label{IV.11.10.9-fr}
Les notations étant celles de (11.10.5), supposons que $X$ soit un
$S$-préschéma localement noethérien et que les $Z_\lambda$ soient tous
$S$-plats. Pour tout $s\in S$, soient
$X_s=X\times_S\operatorname{Spec}(k(s))$,
$(Z_\lambda)_s=Z_\lambda\times_S\operatorname{Spec}(k(s))$,
$(f_\lambda)_s=f_\lambda\times1:(Z_\lambda)_s\to X_s$. Pour que la famille
$(f_\lambda)$ soit universellement schématiquement dominante relativement à
$S$, il faut et il suffit que, pour tout $s\in S$, la famille
$((f_\lambda)_s)$ soit schématiquement dominante.
\end{theorem}

\begin{proposition}[11.10.10]
\label{IV.11.10.10-fr}
Soient $f:X\to S$ un morphisme plat localement de présentation finie, $U$ un
ouvert de $X$. Pour que $U$ soit universellement schématiquement dense dans
$X$ relativement à $S$, il faut et il suffit que, pour tout $s\in S$,
$U_s=U\cap X_s$ soit schématiquement dense dans $X_s$ (ou, ce qui revient au
même, que l'on ait $\operatorname{Ass}(\mathcal O_{X_s})\subset U_s$).
\end{proposition}
% IV3_B25_P172_END
